%% USPSC-modelo.tex
% ---------------------------------------------------------------
% USPSC: Modelo de Trabalho Academico (tese de doutorado, dissertacao de
% mestrado e trabalhos monograficos em geral) em conformidade com 
% ABNT NBR 14724:2011: Informacao e documentacao - Trabalhos academicos -
% Apresentacao
%----------------------------------------------------------------
%% Esta é uma customização do abntex2-modelo-trabalho-academico.tex de v-1.9.5 laurocesar 
%% para as Unidades do Campus USP de São Carlos:
%% EESC - Escola de Engenharia de São Carlos
%% IAU - Instituto de Arquitetura e Urbanismo
%% ICMC - Instituto de Ciências Matemáticas e de Computação
%% IFSC - Instituto de Física de São Carlos
%% IQSC - Instituto de Química de São Carlos
%%
%% Este trabalho utiliza a classe USPSC.cls que é mantida pela seguinte equipe:
%% 
%% Programação:
%%   - Marilza Aparecida Rodrigues Tognetti - marilza@sc.usp.br (PUSP-SC)
%%   - Ana Paula Aparecida Calabrez - aninha@sc.usp.br (PUSP-SC)
%% Normalização e Padronização:
%%   - Brianda de Oliveira Ordonho Sigolo - brianda@usp.br (IAU)
%%   - Elena Luzia Palloni Gonçalves - elena@sc.usp.br (EESC)
%%   - Eliana de Cássia Aquareli Cordeiro - eliana@iqsc.usp.br (IQSC)
%%   - Flávia Helena Cassin - cassinp@sc.usp.br (EESC)
%%   - Maria Cristina Cavarette Dziabas - mcdziaba@ifsc.usp.br (IFSC)
%%   - Regina Célia Vidal Medeiros - rcvmat@icmc.usp.br (ICMC)
%%
%% O USPSC-modelo.tex utiliza:	
%%  USPSC.cls e USPSC1.cls
%% 	USPSC-modelo-references.bib
%%	USPSC-modelo.tex
%%	USPSC-unidades.tex
%%	Um dos arquivos com dados pre-textuais abaixo, em conformidade com a Unidade de vínculo do autor:
%%				USPSC-pre-textual-EESC.tex
%%				USPSC-pre-textual-IAU.tex
%%				USPSC-pre-textual-ICMC.tex
%%				USPSC-pre-textual-IFSC.tex
%%				USPSC-pre-textual-IQSC.tex
%%				USPSC-pre-textual-OUTRO.tex
%%	USPSC-fichacatalografica.tex ou fichacatalografica.pdf
%%	folhadeaprovacao.pdf
%%	USPSC-Cap1-Introducao.tex
%%	USPSC-Cap2-Desenvolvimento.tex
%%	USPSC-Cap3-Citacoes.tex
%%	USPSC-Cap4-referencias.tex
%%	USPSC-Cap5-Conclusao.tex
%%	USPSC-Apendices.tex
%%	USPSC-Anexos.tex
%%	USPSC-AcentuacaoLaTeX.tex
%%	USPSC-LetrasGregas.tex
%%	USPSC-SimbolosUteis.tex

%----------------------------------------------------------------
%% Sobre a classe abntex2.cls:
%% abntex2.cls, v-1.9.5 laurocesar
%% Copyright 2012-2015 by abnTeX2 group at https://www.abntex.net.br/ 
%%
%----------------------------------------------------------------

\documentclass[
% -- opções da classe memoir --
12pt,		% tamanho da fonte
openright,	% capítulos começam em pág ímpar (insere página vazia caso preciso)
twoside,  % para impressão em anverso (frente) e verso. Oposto a oneside - Nota: utilizar \imprimirfolhaderosto*
%oneside, % para impressão em páginas separadas (somente anverso) -  Nota: utilizar \imprimirfolhaderosto
% inclua uma % antes do comando twoside e exclua a % antes do oneside 
a4paper,			% tamanho do papel. 
% -- opções da classe abntex2 --
chapter=TITLE,		% títulos de capítulos convertidos em letras maiúsculas
% -- opções do pacote babel --
english,			% idioma adicional para hifenização
french,				% idioma adicional para hifenização
spanish,			% idioma adicional para hifenização
brazil				% o último idioma é o principal do documento
% {USPSC} configura o cabeçalho contendo apenas o número da página
]{USPSC}
%]{USPSC1}
% Inclua % antes de ]{USPSC} e retire a % antes de %]{USPSC1}
% para utilizar o cabeçalho diferenciado para as páginas pares e ímpares como indicado abaixo:
%- páginas ímpares: cabeçalho com seções ou subseções e o número da página
%- páginas pares: cabeçalho com o número da página e o título do capítulo 
% ---

% ---
% Pacotes básicos - Fundamentais 
% ---
\usepackage[T1]{fontenc}		% Seleção de códigos de fonte.
\usepackage[utf8]{inputenc}		% Codificação do documento (conversão automática dos acentos)
\usepackage{lmodern}			% Usa a fonte Latin Modern
% Para utilizar a fonte Times New Roman, inclua uma % no início do comando acima  "\usepackage{lmodern}"
% Abaixo, tire a % antes do comando  \usepackage{times}
%\usepackage{times}		    	% Usa a fonte Times New Roman	
% Lembre-se de alterar a fonte no comando que imprime o preâmbulo no arquivo da Classe USPSC.cls				
\usepackage{lastpage}			% Usado pela Ficha catalográfica
\usepackage{indentfirst}		% Indenta o primeiro parágrafo de cada seção.
\usepackage{color}				% Controle das cores
\usepackage{graphicx}			% Inclusão de gráficos
\usepackage{float} 				% Fixa tabelas e figuras no local exato
\usepackage{chemfig,chemmacros} % Para escrever reações químicas
\usepackage{microtype} 			% para melhorias de justificação
\usepackage{pdfpages}
\usepackage{makeidx}            % para gerar índice remissimo
% ---

% ---
% Pacotes de citações
% Citações padrão ABNT
% ---
% Sistemas de chamada: autor-data ou numérico.
% Sistema autor-data
\usepackage[alf,abnt-emphasize=bf, abnt-thesis-year=both, abnt-repeated-author-omit=yes, abnt-last-names=abnt, abnt-etal-cite,abnt-etal-list=3, abnt-etal-text=default, abnt-and-type=e, abnt-doi=doi, abnt-url-package=none, abnt-verbatim-entry=no]{abntex2cite}

% Para o IQSC, que indica todos os autores nas referências, incluir % no início do comando acima e retirar a % do comando abaixo 

%\usepackage[alf,abnt-emphasize=bf, abnt-thesis-year=both, abnt-repeated-author-omit=yes, abnt-last-names=abnt, abnt-etal-cite,abnt-etal-list=0, abnt-etal-text=default, abnt-and-type=e]{abntex2cite}

% Sistema Numérico
% Para citações numéricas, sistema adotado pelo IFSC, incluir % no início do comando acima e retirar a % do comando abaixo 
%\usepackage[num,overcite,abnt-emphasize=bf, abnt-thesis-year=both, abnt-repeated-author-omit=yes, abnt-last-names=abnt, abnt-etal-cite,abnt-etal-list=0, abnt-etal-text=default, abnt-and-type=e]{abntex2cite}

% Complementarmente, verifique as instruções abaixo sobre os Pacotes de Nota de rodapé
% ---
% Pacotes de Nota de rodapé
% Configurações de nota de rodapé

% O presente modelo adota o formato numérico para as notas de rodapés quando utiliza o sistema de chamada autor-data para citações e referências. Para utilizar o sistema de chamada numérico para citações e referências, habilitar um dos comandos abaixo.
% Há diversa opções para nota de rodapé no Sistema Numérico.  Para o IFSC, habilitade o comando abaixo.

%\renewcommand{\thefootnote}{\fnsymbol{footnote}}  %Comando para inserção de símbolos em nota de rodapé

% Outras opções para nota de rodapé no Sistema Numérico:
%\renewcommand{\thefootnote}{\alph{footnote}}      %Comando para inserção de letras minúscula em nota de rodapé
%\renewcommand{\thefootnote}{\Alph{footnote}}      %Comando para inserção de letras maiúscula em nota de rodapé
%\renewcommand{\thefootnote}{\roman{footnote}}     %Comando para inserção de números romanos minúsculos  em nota de rodapé
%\renewcommand{\thefootnote}{\Roman{footnote}}     %Comando para inserção de números romanos minúsculos  em nota de rodapé

\renewcommand{\footnotesize}{\small} %Comando para diminuir a fonte das notas de rodapé

 % ---
 % Pacote para agrupar a citação numérica consecutiva
 % Quando for adotado o Sistema Numérico, a exemplo do IFSC, habilite 
 % o pacote cite abaixo retirando a porcentagem antes do comando abaixo
 %\usepackage[superscript]{cite}	

% ---
% Pacotes adicionais, usados apenas no âmbito do Modelo Canônico do abnteX2
% ---
\usepackage{lipsum}				% para geração de dummy text
% ---




% pacotes de tabelas
\usepackage{multicol}	% Suporte a mesclagens em colunas
\usepackage{multirow}	% Suporte a mesclagens em linhas
\usepackage{longtable}	% Tabelas com várias páginas
\usepackage{threeparttablex}    % notas no longtable
\usepackage{array}

%---
% Configurações para o pacote chemfig
%\chemsetup[chemformula]{format=\sffamily}
\renewcommand*\printatom[1]{\ensuremath{\mathsf{#1}}}
\setatomsep{2em}
\setdoublesep{.6ex}
\setbondstyle{semithick}
%---
% Configurando o ambiente para utilizar os recursos de frases pre-prontas do mhchem
\newenvironment{rslist}%
{%
	\begin{labeling}% environment from KOMA-script
		{\rsnumber{R39/23/24/25}}% R39/23/24/25 is longest label
	}{%
\end{labeling}%
}%
% Definição de comando para utilizar os recursos de frases pre-prontas do mhchem
\newcommand{\rs}[2][]{\item[{\rsnumber[#1]{#2}}] \rsphrase{bb}}
% ---

% ---
% DADOS INICIAIS - Define sigla com título, área de concentração e opção do Programa 
% Consulte a tabela referente aos Programas, áreas e opções de sua unidade contante do
% arquivo USPSC-Siglas estabelecidas para os Programas de Pós-Graduação por Unidade.xlsx 
% ou nos APÊNDICES A-E
\siglaunidade{EESC}
\programa{DCEA}
% Os demais dados deverão ser fornecidos no arquivo USPSC-pre-textual-UUUU, onde UUUU é a sigla da Unidade. 
% Exemplo:USPSC-pre-textual-IFSC.tex
% ---
% Configurações de aparência do PDF final
% alterando o aspecto da cor azul
\definecolor{blue}{RGB}{41,5,195}

% informações do PDF
\makeatletter
\hypersetup{
	%pagebackref=true,
	pdftitle={\@title}, 
	pdfauthor={\@author},
	pdfsubject={\imprimirpreambulo},
	pdfcreator={LaTeX with abnTeX2},
	pdfkeywords={abnt}{latex}{abntex}{USPSC}{trabalho acadêmico}, 
	colorlinks=true,       		% false: boxed links; true: colored links
	linkcolor=blue,          	% color of internal links
	citecolor=blue,        		% color of links to bibliography
	filecolor=magenta,      		% color of file links
	urlcolor=blue,
	bookmarksdepth=4
}
\makeatother
% --- 
% --- 
% Espaçamentos entre linhas e parágrafos 
% --- 

% O tamanho do parágrafo é dado por:
\setlength{\parindent}{1.3cm}

% Controle do espaçamento entre um parágrafo e outro:
\setlength{\parskip}{0.2cm}  % tente também \onelineskip

% ---
% compila o sumário e índice
\makeindex
% ---

% ----
% Início do documento
% ----
\begin{document}

% Seleciona o idioma do documento (conforme pacotes do babel)
\selectlanguage{brazil}
% Se o idioma do texto for inglês, inclua uma % antes do 
%      comando \selectlanguage{brazil} e 
%      retire a % antes do comando abaixo
%\selectlanguage{english}

% Retira espaço extra obsoleto entre as frases.
\frenchspacing 

% --- Formatação dos Títulos
\renewcommand{\ABNTEXchapterfontsize}{\fontsize{12}{12}\bfseries}
\renewcommand{\ABNTEXsectionfontsize}{\fontsize{12}{12}\bfseries}
\renewcommand{\ABNTEXsubsectionfontsize}{\fontsize{12}{12}\normalfont}
\renewcommand{\ABNTEXsubsubsectionfontsize}{\fontsize{12}{12}\normalfont}
\renewcommand{\ABNTEXsubsubsubsectionfontsize}{\fontsize{12}{12}\normalfont}


% ----------------------------------------------------------
% ELEMENTOS PRÉ-TEXTUAIS
% ----------------------------------------------------------
% ---
% Capa
% ---
\imprimircapa
% ---
% Folha de rosto
% (o * indica impressão em anverso (frente) e verso )
% ---
\imprimirfolhaderosto*
%\imprimirfolhaderosto
% ---

% ---
% Inserir a ficha catalográfica em pdf
% ---
% A biblioteca da sua Unidade lhe fornecerá um PDF com a ficha
% catalográfica definitiva. 
% Quando estiver com o documento, salve-o como PDF no diretório
% do seu projeto como fichacatalografica.pdf e iclua o arquivo
% utilizando o comando abaixo:
\begin{fichacatalografica}
   \includepdf{fichacatalografica.pdf}
\end{fichacatalografica}
% Se você optar por elaborar a ficha catalográfica, deverá 
% incluir uma % antes das 3 linhas acima e tirar a % antes
% do comando % ---
% Inserir a ficha bibliografica
% ---
% Isto é um exemplo de Ficha Catalográfica, ou ``Dados internacionais de
% catalogação-na-publicação''. Você pode utilizar este modelo como referência. 
% Porém, provavelmente a biblioteca da sua universidade lhe fornecerá um PDF
% com a ficha catalográfica definitiva após a defesa do trabalho. Quando estiver
% com o documento, salve-o como PDF no diretório do seu projeto e substitua todo
% o conteúdo de implementação deste arquivo pelo comando abaixo:
%
\begin{fichacatalografica}
	\hspace{-1.4cm}
	\imprimirnotaautorizacao \\ \\
	%\sffamily
	\vspace*{\fill}					% Posição vertical
	\begin{center}					% Minipage Centralizado
		\imprimirnotabib \\
\begin{table}[htb]
	\scriptsize
	\centering	
	\begin{tabular}{|p{0.9cm} p{8.7cm}|}
		\hline
	      & \\
		  &	  \imprimirautorficha     \\
		
		 \imprimircutter & 
							\hspace{0.4cm}\imprimirtitulo~  / ~\imprimirautor~ ;  ~\imprimirorientadorcorpoficha. -- 	\imprimirlocal, \imprimirdata.   \\
		
		  &  % Para incluir nota referente à versão corrigida no corpo da ficha,
			  % incluir % no início da linha acima e tirar a % do início da linha abaixo
			  %	\hspace{0.4cm} \imprimirtitulo~  / ~\imprimirautor~ ; ~\imprimirorientadorcorpoficha~- ~\imprimirnotafolharosto. -- \imprimirlocal, \imprimirdata.  \\
		
			\hspace{0.4cm}\pageref{LastPage} p. : il. (algumas color.) ; 30 cm.\\ 
		  & \\
		  & 
		    \hspace{0.4cm}\imprimirnotaficha ~--~ 
						  \imprimirunidademin, 
						  \imprimiruniversidademin, 
		                  \imprimirdata. \\ 
		  & \\                 
		   % Para incluir nota referente à versão corrigida em notas,
		    % incluir uma % no início da linha acima e	
		    % tirar a % do início da linha abaixo
		    % & \hspace{0.4cm}\imprimirnotafolharosto \\ 
		  & \\ 
		  & \hspace{0.4cm}1. LaTeX. 2. abnTeX. 3. Classe USPSC. 4. Editoração de texto. 5. Normalização da documentação. 6. Tese. 7. Dissertação. 8. Documentos (elaboração). 9. Documentos eletrônicos. I. \imprimirorientadorficha. 
		   II. Título. \\
	
		     %Se houver co-orientador, inclua % antes da linha (antes de II. Título.) 
		     %          e tire a % antes do comando abaixo 
		     %III. Título. \\   
		  \hline
	\end{tabular}
\end{table}
	\end{center}
\end{fichacatalografica}
% ---


%% ---
% Inserir a ficha bibliografica
% ---
% Isto é um exemplo de Ficha Catalográfica, ou ``Dados internacionais de
% catalogação-na-publicação''. Você pode utilizar este modelo como referência. 
% Porém, provavelmente a biblioteca da sua universidade lhe fornecerá um PDF
% com a ficha catalográfica definitiva após a defesa do trabalho. Quando estiver
% com o documento, salve-o como PDF no diretório do seu projeto e substitua todo
% o conteúdo de implementação deste arquivo pelo comando abaixo:
%
\begin{fichacatalografica}
	\hspace{-1.4cm}
	\imprimirnotaautorizacao \\ \\
	%\sffamily
	\vspace*{\fill}					% Posição vertical
	\begin{center}					% Minipage Centralizado
		\imprimirnotabib \\
\begin{table}[htb]
	\scriptsize
	\centering	
	\begin{tabular}{|p{0.9cm} p{8.7cm}|}
		\hline
	      & \\
		  &	  \imprimirautorficha     \\
		
		 \imprimircutter & 
							\hspace{0.4cm}\imprimirtitulo~  / ~\imprimirautor~ ;  ~\imprimirorientadorcorpoficha. -- 	\imprimirlocal, \imprimirdata.   \\
		
		  &  % Para incluir nota referente à versão corrigida no corpo da ficha,
			  % incluir % no início da linha acima e tirar a % do início da linha abaixo
			  %	\hspace{0.4cm} \imprimirtitulo~  / ~\imprimirautor~ ; ~\imprimirorientadorcorpoficha~- ~\imprimirnotafolharosto. -- \imprimirlocal, \imprimirdata.  \\
		
			\hspace{0.4cm}\pageref{LastPage} p. : il. (algumas color.) ; 30 cm.\\ 
		  & \\
		  & 
		    \hspace{0.4cm}\imprimirnotaficha ~--~ 
						  \imprimirunidademin, 
						  \imprimiruniversidademin, 
		                  \imprimirdata. \\ 
		  & \\                 
		   % Para incluir nota referente à versão corrigida em notas,
		    % incluir uma % no início da linha acima e	
		    % tirar a % do início da linha abaixo
		    % & \hspace{0.4cm}\imprimirnotafolharosto \\ 
		  & \\ 
		  & \hspace{0.4cm}1. LaTeX. 2. abnTeX. 3. Classe USPSC. 4. Editoração de texto. 5. Normalização da documentação. 6. Tese. 7. Dissertação. 8. Documentos (elaboração). 9. Documentos eletrônicos. I. \imprimirorientadorficha. 
		   II. Título. \\
	
		     %Se houver co-orientador, inclua % antes da linha (antes de II. Título.) 
		     %          e tire a % antes do comando abaixo 
		     %III. Título. \\   
		  \hline
	\end{tabular}
\end{table}
	\end{center}
\end{fichacatalografica}
% ---


% As informações que compõem a ficha catalográfica estão 
% definidos no arquivo USPSC-pre-textual-UUUU.tex
% ---


% ---
% Inserir errata
% ---

\begin{errata}
	\OnehalfSpacing 			
	A errata é um elemento opcional, que consiste de uma lista de erros da obra, precedidos pelas folhas e linhas onde eles ocorrem e seguidos pelas correções correspondentes. Deve ser inserida logo após a folha de rosto e conter a referência do trabalho para facilitar sua identificação, conforme a ABNT NBR 14724 \cite{nbr14724}.
	
	Modelo de Errata:
		
	\begin{flushleft} 
			\setlength{\absparsep}{0pt} % ajusta o espaçamento da referência	
			\SingleSpacing 
			\imprimirautorabr~ ~\textbf{\imprimirtitulo}.	\imprimirdata. \pageref{LastPage}p. 
			%Substitua p. por f. quando utilizar oneside em \documentclass
			%\pageref{LastPage}f.
			\imprimirtipotrabalho~-~\imprimirinstituicao, \imprimirlocal, \imprimirdata. 
 	\end{flushleft}
\vspace{\onelineskip}
\OnehalfSpacing 
\center
\textbf{ERRATA}
\vspace{\onelineskip}
\OnehalfSpacing 
\begin{table}[htb]
	\center
	\footnotesize
	\begin{tabular}{p{1.4cm} p{1cm} p{3cm} p{3cm} }
		\hline
		\textbf{Folha} & \textbf{Linha}  & \textbf{Onde se lê}  & \textbf{Leia-se}  \\
			\hline
			1 & 10 & auto-conclavo & autoconclavo\\
		\hline
	\end{tabular}
\end{table}


\end{errata}
% ---

% ---
% Inserir folha de aprovação
% ---

% A Folha de aprovação é um elemento obrigatório da NBR 4724/2011 (seção 4.2.1.3). 
% Após a defesa/aprovação do trabalho, gere o arquivo folhadeaprovacao.pdf da página assinada pela banca 
% e iclua o arquivo utilizando o comando abaixo:
\includepdf{folhadeaprovacao.pdf}

% Alternativa para a Folha de Aprovação:
% Se for a sua opção elaborar uma folha de aprovação, insira uma % antes do comando acima que inclui o arquivo folhadeaprovacao.pdf,
% tire o % do comando abaixo e altere o arquivo folhadeaprovacao.tex conforme suas necessidades
%% Alternativa para a Folha de Aprova��o 
% Se esta for a sua op��o, exclua inclus�o feita acima do arquivo folhadeaprovacao.pdf
%
\begin{folhadeaprovacao}
  \begin{center}
       {\ABNTEXchapterfont\bfseries\large\imprimirautor}
	 \vspace*{2cm}
   
    \begin{center}
      \ABNTEXchapterfont\bfseries\Large\imprimirtitulo
    \end{center}
		\vspace*{2cm}
		\hspace{.45\textwidth}
    \begin{minipage}{.5\textwidth}
        \imprimirpreambulo
    \end{minipage}
		\vspace*{2cm}
    %\vspace*{\fill}
	\end{center}

  \begin{center}
	  {\ABNTEXchapterfont\bfseries\large\ {Data de defesa: 02 de outubro de 2015} \\}
		\vspace*{\fill}
	  {\ABNTEXchapterfont\bfseries\large\ {Comiss\~ao Julgadora:} \\}
		%Trabalho aprovado. \imprimirlocal, 2 de outubro de 2015:
		
		%\assinatura{\textbf{\imprimirorientador} \\ Orientador} 
		\assinatura{\textbf{\imprimirorientador} \\}
    % Se for ORIENTADOR, inclua % no in�cio do comando abaixo e tire a % do pr�ximo comando 
		\renewcommand{\orientadorname}{Orientadora}
		%\renewcommand{\orientadorname}{Orientador}

		\imprimirorientadorRotulo
		\par
		\assinatura{\textbf{Professor} \\ Convidado1}
	
		\assinatura{\textbf{Professor} \\ Convidado2}
		%\assinatura{\textbf{Professor} \\ Convidado3}
		%\assinatura{\textbf{Professor} \\ Convidado4}
		%\begin{center}
		\vspace*{\fill}
   	{\ABNTEXchapterfont\bfseries\large\imprimirlocal}
    \par
    {\ABNTEXchapterfont\bfseries\large\imprimirdata}
\end{center}
\end{folhadeaprovacao}
% ---

\includepdf{PaginaEmBranco.pdf}

% ---
% Dedicatória
% ---
\begin{dedicatoria}
   \vspace*{\fill}
   \centering
   \noindent
   \textit{ Este trabalho é dedicado aos alunos da USP, como uma contribuição\\
  das Bibliotecas do Campus USP de São Carlos para o desenvolvimento\\
	e disseminação da pesquisa científica da Universidade.} \vspace*{\fill}
\end{dedicatoria}
% ---

% ---
% Agradecimentos
% ---=====
\begin{agradecimentos}
A motivação para o desenvolvimento da classe USPSC e dos modelos de dissertações e teses foi decorrente de solicitações de usuários das Bibliotecas do Campus USP de São Carlos. 

O Grupo desenvolvedor do Pacote USPSC, atualmente composto da Classe USPSC e do  Modelo para teses e dissertações em \LaTeX\ utilizando a classe USPSC, agradece especialmente ao Luis Olmes, doutorando do Instituto de Ciências Matemáticas e de Computação (ICMC) da Universidade de São Paulo (USP), pelas primeiras orientações sobre o \LaTeX\ .

Agradecemos ao Lauro César Araujo pelo desenvolvimento da classe  \abnTeX, modelos canônicos e tantas outras contribuições que nos permitiu o desenvolvimento da classe USPSC e seus modelos.

Os nossos agradecimentos aos integrantes do primeiro
projeto abn\TeX\, Gerald Weber, Miguel Frasson, Leslie H. Watter, Bruno Parente Lima, Flávio de Vasconcellos Corrêa, Otavio Real
Salvador, Renato Machnievscz, e a todos que contribuíram para que a produção de trabalhos acadêmicos em conformidade com
as normas ABNT com \LaTeX\ fosse possível.

Agradecemos ao grupo de usuários
\emph{latex-br}{\url{http://groups.google.com/group/latex-br}}, aos integrantes do grupo
\emph{\abnTeX}{\url{http://groups.google.com/group/abntex2}  e \url{http://www.abntex.net.br/}}~que contribuem para a evolução do \abnTeX.
\end{agradecimentos}
% ---

% ---
% Epígrafe
% ---
\begin{epigrafe}
    \vspace*{\fill}
	\begin{flushright}
		\textit{``O estudo, a busca da verdade e da beleza são domínios \\
		em que nos é consentido sermos crianças por toda a vida.''\\
		Albert Einstein}
	\end{flushright}
\end{epigrafe}
% ---

% ---
% Resumo
% ---
% Se o idioma do texto for em inglês, o abstract deve preceder o resumo
% resumo em português
\setlength{\absparsep}{18pt} % ajusta o espaçamento dos parágrafos do resumo		
\begin{resumo}
	\begin{flushleft} 
			\setlength{\absparsep}{0pt} % ajusta o espaçamento da referência	
			\SingleSpacing 
			\imprimirautorabr~ ~\textbf{\imprimirtitulo}.	\imprimirdata. \pageref{LastPage}p. 
			%Substitua p. por f. quando utilizar oneside em \documentclass
			%\pageref{LastPage}f.
			\imprimirtipotrabalho~-~\imprimirinstituicao, \imprimirlocal, \imprimirdata. 
 	\end{flushleft}
\OnehalfSpacing 			
 O resumo deve ressaltar o  objetivo, o método, os resultados e as conclusões do documento. A ordem e a extensão  destes itens dependem do tipo de resumo (informativo ou indicativo) e do  tratamento que cada item recebe no documento original. O resumo deve ser
 precedido da referência do documento, com exceção do resumo inserido no
 próprio documento. (\ldots)  Salientamos que os Programas de Pós-Graduação de algumas Unidades exigem o titulo da dissertação ou tese em inglês, tornando necessário a inclusão das referências nos resumos e abstracts, o que foi adotado no presente modelo. As palavras-chave devem figurar logo abaixo do  resumo, antecedidas da expressão Palavras-chave:, separadas entre si por  ponto e finalizadas também por ponto \cite{nbr6028}.
 

 \textbf{Palavras-chave}: LaTeX. abnTeX. Classe USPSC. Editoração de texto. Normalização da documentação. Tese. Dissertação. Documentos (elaboração). Documentos eletrônicos. 
\end{resumo}

% ---
% Abstract
% ---
\autor{Silva, M. J.}
\begin{resumo}[Abstract]
 \begin{otherlanguage*}{english}
	\begin{flushleft} 
		\setlength{\absparsep}{0pt} % ajusta o espaçamento dos parágrafos do resumo		
 		\SingleSpacing 
 		\imprimirautorabr~ ~\textbf{\imprimirtitleabstract}.	\imprimirdata.  \pageref{LastPage}p. 
		%Substitua p. por f. quando utilizar oneside em \documentclass
		%\pageref{LastPage}f.
		\imprimirtipotrabalho~-~\imprimirinstituicao, \imprimirlocal, 	\imprimirdata. 
 	\end{flushleft}
	\OnehalfSpacing 
   This is the english abstract.

   \vspace{\onelineskip}
 
   \noindent 
   \textbf{Keywords}: LaTeX. abnTeX. USPSC class. Text editoration. Standardization of documentation. Thesis. Dissertation. Documents (development). Electronic documents.
 \end{otherlanguage*}
\end{resumo}

% ---
% inserir lista de ilustrações
% ---
\pdfbookmark[0]{\listfigurename}{lof}
\listoffigures*
\cleardoublepage
% ---

% ---
% inserir lista de tabelas
% ---
\pdfbookmark[0]{\listtablename}{lot}
\listoftables*
\cleardoublepage
% ---

% ---
% inserir lista de quadros
% ---
\pdfbookmark[0]{\listofquadroname}{loq}
\listofquadro*
\cleardoublepage
% ---

% ---
% inserir lista de abreviaturas e siglas
% ---
\begin{siglas}
    \item[ABNT] Associação Brasileira de Normas Técnicas
    \item[abnTeX] ABsurdas Normas para TeX
	\item[EESC] Escola de Engenharia de São Carlos
	\item[IAU] Instituto de Arquitetura e Urbanismo
	\item[IBGE] Instituto Brasileiro de Geografia e Estatística
	\item[ICMC] Instituto de Ciências Matemáticas e de Computação
	\item[IFSC] Instituto de Física de São Carlos
	\item[IQSC] Instituto de Química de São Carlos
	\item[USP] Universidade de São Paulo
	\item[USPSC] Campus USP de São Carlos
	
\end{siglas}
% ---

% ---
% inserir lista de símbolos
% ---
\begin{simbolos}
  \item[$ \Gamma $] Letra grega Gama
  \item[$ \Lambda $] Lambda
  \item[$ \zeta $] Letra grega minúscula zeta
  \item[$ \in $] Pertence
\end{simbolos}
% ---
% ---
% inserir o sumario
% ---
\pdfbookmark[0]{\contentsname}{toc}
\tableofcontents*
\cleardoublepage
% ---
% ----------------------------------------------------------
% ELEMENTOS TEXTUAIS
% ----------------------------------------------------------
\textual
% Os capítulos são inseridos como arquivos externos 

% Capítulo 1 - Introdução
% ---
\include{USPSC-Cap1-Introducao}
% ---

% ---
% Capítulo 2
% ---
\include{USPSC-Cap2-Desenvolvimento}
% ---
% Capítulo 3 - Citações
% ---
\include{USPSC-Cap3-Citacoes}

% ---
% Capítulo 4 - Referencias
% ---
\include{USPSC-Cap4-Referencias}
% ---

% Capítulo 5 - Conclusão
% ---
\include{USPSC-Cap5-Conclusao}
% ---

% ----------------------------------------------------------
% ELEMENTOS PÓS-TEXTUAIS
% ----------------------------------------------------------
\postextual
% ----------------------------------------------------------

% -----------------------------------------------------------
% Referências bibliográficas
% ----------------------------------------------------------
\bibliography{USPSC-modelo-references}


% ----------------------------------------------------------
% Glossário
% ----------------------------------------------------------
%
% Consulte o manual da classe abntex2 para orientações sobre o glossário.
%
%\glossary

% ----------------------------------------------------------
% Apêndices
% ----------------------------------------------------------
%% USPSC-Apendice.tex
% ---
% Inicia os apêndices
% ---

\begin{apendicesenv}
% Imprime uma página indicando o início dos apêndices
\partapendices
\chapter{Apendice(s)}
Elemento opcional, que consiste em texto ou documento elaborado pelo autor, a fim de complementar sua argumentação, conforme a ABNT NBR 14724 \cite{nbr14724}.

Os apêndices devem ser identificados por letras maiúsculas consecutivas, seguidas de hífen e pelos respectivos títulos. Excepcionalmente, utilizam-se letras maiúsculas dobradas na identificação dos apêndices, quando esgotadas as 26 letras do alfabeto. A paginação deve ser contínua, dando seguimento ao texto principal. \cite{sibi2009}.

\chapter{Siglas dos Programas de Pós-Graduação da EESC}
\index{quadros}O \autoref{quadro-eesc} relaciona as siglas estabelecidas para os programas de pós-graduação da EESC.

\begin{quadro}[htb]
\ABNTEXfontereduzida
%\caption[Siglas dos Programas de Pós-Graduação da EESC]{Siglas dos Programas de Pós-Graduação da EESC]{Siglas dos Programas de Pós-Graduação da EESC}
\caption[Siglas dos Programas de Pós-Graduação da EESC]{Siglas dos Programas de Pós-Graduação da EESC} 
\label{quadro-eesc}
\begin{tabular}{|p{6.0cm}|p{4.5cm}|p{2.0cm}|p{1.75cm}|}
%\multicolumn{4}{c}%
%{{\tablename\ \thetable{} -- Siglas dos Programas de Pós-Graduação da EESC}} \\
\multicolumn{4}{r}{{(continua)}} \\ 
  \hline
   \textbf{PROGRAMA} & \textbf{ÁREA DE CONCENTRAÇÃO} & \textbf{TÍTULO} & \textbf{SIGLA}  \\
    \hline
Programa de Pós-Graduação em Ciências da Engenharia Ambiental & Ciências da Engenharia Ambiental & Doutor & DCEA \\
Programa de Pós-Graduação em Ciências da Engenharia Ambiental & Ciências da Engenharia Ambiental & Mestre & MCEA \\
Programa de Pós-Graduação em Ciência e Engenharia de Materiais & Desenvolvimento, Caracterização e Aplicação de Materiais  & Doutor & DCEM \\
Programa de Pós-Graduação em Ciência e Engenharia de Materiais & Desenvolvimento, Caracterização e Aplicação de Materiais & Mestre & MCEM \\
Programa de Pós-Graduação em Engenharia Civil (Engenharia de Estruturas) & Estruturas & Doutor & DEE \\
Programa de Pós-Graduação em Engenharia Civil (Engenharia de Estruturas) & Estruturas & Mestre & MEE \\
Programa de Pós-Graduação em Engenharia de Produção & Economia, Organizações e Gestão do Conhecimento & Doutor & DEPE \\
Programa de Pós-Graduação em Engenharia de Produção & Economia, Organizações e Gestão do Conhecimento & Mestre & MEPE \\
Programa de Pós-Graduação em Engenharia de Produção & Processos e Gestão de Operações & Doutor & DEPP \\
Programa de Pós-Graduação em Engenharia de Produção & Processos e Gestão de Operações & Mestre & MEPP \\
Programa de Pós-Graduação em Engenharia de Transportes & Infraestrutura de Transportes & Doutor & DETI \\
Programa de Pós-Graduação em Engenharia de Transportes & Infraestrutura de Transportes & Mestre & METI \\
Programa de Pós-Graduação em Engenharia de Transportes & Planejamento e Operação de Sistemas de Transporte & Doutor & DETP \\
Programa de Pós-Graduação em Engenharia de Transportes & Planejamento e Operação de Sistemas de Transporte & Mestre & METP \\
Programa de Pós-Graduação em Engenharia de Transportes & Transportes & Doutor & DETT \\
Programa de Pós-Graduação em Engenharia de Transportes & Transportes & Mestre & METT \\
Programa de Pós-Graduação em Engenharia Elétrica & Processamento de Sinais e Intrumentação & Doutor & DEEP \\

\end{tabular}
\end{quadro} 

% o comando \clearpage é necessário para deixar o final da tabela o topo da página, sem ele o final da tabela é centralizado verticalmente na página 
\clearpage
\begin{quadro}[htb]
	\ABNTEXfontereduzida
\begin{tabular}{|p{6.0cm}|p{4.5cm}|p{2.0cm}|p{1.75cm}|}	
 \multicolumn{4}{c}%
	{{\quadroname\ \thequadro{} -- Siglas dos Programas de Pós-Graduação da EESC}} \\
	\multicolumn{4}{r}{{(continuação)}} \\
	 \hline
   \textbf{PROGRAMA} & \textbf{ÁREA DE CONCENTRAÇÃO} & \textbf{TÍTULO} & \textbf{SIGLA}  \\
		 \hline
Programa de Pós-Graduação em Engenharia Elétrica & Processamento de Sinais e Intrumentação & Mestre & MEEP \\
Programa de Pós-Graduação em Engenharia Elétrica & Sistemas Dinâmicos & Doutor & DEED \\
Programa de Pós-Graduação em Engenharia Elétrica & Sistemas Dinâmicos & Mestre & MEED \\
Programa de Pós-Graduação em Engenharia Elétrica & Sistemas Elétricos de Potência & Doutor & DEEE \\
Programa de Pós-Graduação em Engenharia Elétrica & Sistemas Elétricos de Potência & Mestre & MEEE \\
Programa de Pós-Graduação em Engenharia Elétrica & Telecomunicações & Doutor & DEET \\
Programa de Pós-Graduação em Engenharia Elétrica & Telecomunicações & Mestre & MEET \\
Programa de Pós-Graduação em Engenharia Hidráulica e Saneamento & Hidráulica e Saneamento & Doutor & DEHS \\
Programa de Pós-Graduação em Engenharia Hidráulica e Saneamento & Hidráulica e Saneamento & Mestre & MEHS \\
Programa de Pós-Graduação em Engenharia Mecânica & Aeronaves & Doutor & DEMA \\
Programa de Pós-Graduação em Engenharia Mecânica & Aeronaves & Mestre & MEMA \\
Programa de Pós-Graduação em Engenharia Mecânica & Dinâmica das Máquinas e Sistemas & Doutor & DEMD \\
Programa de Pós-Graduação em Engenharia Mecânica & Dinâmica das Máquinas e Sistemas & Mestre & MEMD \\
Programa de Pós-Graduação em Engenharia Mecânica & Manufatura & Doutor & DEMF \\
Programa de Pós-Graduação em Engenharia Mecânica & Manufatura & Mestre & MEMF \\
Programa de Pós-Graduação em Engenharia Mecânica & Materiais & Doutor & DEMT \\
Programa de Pós-Graduação em Engenharia Mecânica & Materiais & Mestre & MEMT \\
Programa de Pós-Graduação em Engenharia Mecânica & Projeto Mecânico & Doutor & DEMP \\
Programa de Pós-Graduação em Engenharia Mecânica & Projeto Mecânico & Mestre & MEMP \\
Programa de Pós-Graduação em Engenharia Mecânica & Térmica e Fluídos & Doutor & DEML \\
Programa de Pós-Graduação em Engenharia Mecânica & Térmica e Fluídos & Mestre & MEML \\
Programa de Pós-Graduação em Geotecnia & Geotecnia & Doutor & DGEO \\
Programa de Pós-Graduação em Geotecnia & Geotecnia & Mestre & MGEO \\
    
\end{tabular}
\end{quadro}

% o comando \clearpage é necessário para deixar o final da tabela o topo da página, sem ele o final da tabela é centralizado verticalmente na página 
\clearpage
\begin{quadro}[htb]
	\ABNTEXfontereduzida
\begin{tabular}{|p{6.0cm}|p{4.5cm}|p{2.0cm}|p{1.75cm}|}
\multicolumn{4}{c}%
	{{\quadroname\ \thequadro{} -- Siglas dos Programas de Pós-Graduação da EESC}} \\
	\multicolumn{4}{r}{{(conclusão)}} \\
\hline
\textbf{PROGRAMA} & \textbf{ÁREA DE CONCENTRAÇÃO} & \textbf{TÍTULO} & \textbf{SIGLA}  \\
\hline    
Programa de Pós-Graduação Interunidades em Bioengenharia & Bioengenharia & Doutor & DIUB \\
Programa de Pós-Graduação Interunidades em Bioengenharia & Bioengenharia & Mestre & MIUB \\
Programa de Pós-Graduação em Rede Nacional para Ensino das Ciências Ambientais & Ensino de Ciências Ambientais & Mestre & MRNECA \\    
    
    \hline
\end{tabular}
\begin{flushleft}
		Fonte: Elaborado pelos autores.\
\end{flushleft}
\end{quadro}

% ----------------------------------------------------------
\chapter{Siglas dos Programas de Pós-Graduação do IAU}
\index{quadros}O \autoref{quadro-iau} relaciona as siglas estabelecidas para os programas de pós-graduação do IAU.
\begin{quadro}[htb]
\ABNTEXfontereduzida
\caption[Siglas dos Programas de Pós-Graduação do IAU]{Siglas dos Programas de Pós-Graduação do IAU}
\label{quadro-iau}
\begin{tabular}{|p{3.5cm}|p{3.5cm}|p{3.5cm}|p{1.5cm}|p{2.25cm}|}
  \hline
   \textbf{PROGRAMA} & \textbf{ÁREA DE CONCENTRAÇÃO} & \textbf{OPÇÃO} & \textbf{TÍTULO} & \textbf{SIGLA}  \\
    \hline
Programa de Pós-Graduação em Arquitetura e Urbanismo & Arquitetura, Urbanismo e Tecnologia &  & Doutor & DAUT\\
Programa de Pós-Graduação em Arquitetura e Urbanismo & Arquitetura, Urbanismo e Tecnologia &  & Mestre & MAUT\\
Programa de Pós-Graduação em Arquitetura e Urbanismo & Teoria e História da Arquitetura e do Urbanismo &  & Doutor & DAUH\\
Programa de Pós-Graduação em Arquitetura e Urbanismo & Teoria e História da Arquitetura e do Urbanismo &  & Mestre & MAUH\\
    \hline

\end{tabular}
\begin{flushleft}
		Fonte: Elaborado pelos autores.\
\end{flushleft}
\end{quadro}

% ----------------------------------------------------------
\chapter{Siglas dos Programas de Pós-Graduação do ICMC}
\index{quadros}O \autoref{quadro-icmc} relaciona as siglas estabelecidas para os programas de pós-graduação do ICMC.
\begin{quadro}[htb]
\ABNTEXfontereduzida
\caption[Siglas dos Programas de Pós-Graduação do ICMC]{Siglas dos Programas de Pós-Graduação do ICMC}
\label{quadro-icmc}
\begin{tabular}{|p{3.5cm}|p{3.5cm}|p{3.5cm}|p{1.5cm}|p{2.25cm}|}
  \multicolumn{5}{r}{{(continua)}} \\ 
  \hline
   \textbf{PROGRAMA} & \textbf{ÁREA DE CONCENTRAÇÃO} & \textbf{OPÇÃO} & \textbf{TÍTULO} & \textbf{SIGLA}  \\
    \hline
		Ciências de Computação e Matemática Computacional	& Ciências de Computação e Matemática Computacional	&   &	Doutor	 & DCCp\\
    Ciências de Computação e Matemática Computacional	& Ciências de Computação e Matemática Computacional	&   &	Mestre	& MCCp\\
		Computer Science and Computational Mathematics & Computer Science and Computational Mathematics	&   &	Doctorate & DCCe\\
		Doctorate Program in Mathematics & Mathematics &   &	Doctorate & DMAe\\
		Interinstitucional de Pós-Graduação em Estatística & Estatística &  & Doutor	 & DESp\\
		Interinstitucional de Pós-Graduação em Estatística & Estatística &  & Mestre & MESp\\
		Join Graduate Program in Statistics & Computer Science and Computational Mathematics &  & Master & MCCe\\
		Join Graduate Program in Statistics & Statistics &  & Doctorate & 	DESe\\
    Join Graduate Program in Statistics & Statistics &  & Master & MESe\\
		Master Program in Mathematics &	Mathematics &  & Master &	MMAe\\
		Mathematics Professional Master\'{}s Program &	Mathematics &	 & Master &	MPMe\\
		Programa de Mestrado Profissional em Matemática & Matemática &  & Mestre & MPMp\\
		\end{tabular}
\end{quadro}

% o comando \clearpage é necessário para deixar o final da tabela o topo da página, sem ele o final da tabela é centralizado verticalmente na página 
\clearpage
\begin{quadro}[htb]
\ABNTEXfontereduzida
\begin{tabular}{|p{3.5cm}|p{3.5cm}|p{3.5cm}|p{1.5cm}|p{2.25cm}|}
	\multicolumn{5}{c}%
	{{\quadroname\ \thequadro{} -- Siglas dos Programas de Pós-Graduação do ICMC}} \\
	\multicolumn{5}{r}{{(conclusão)}} \\
	\hline
   \textbf{PROGRAMA} & \textbf{ÁREA DE CONCENTRAÇÃO} & \textbf{OPÇÃO} & \textbf{TÍTULO} & \textbf{SIGLA}  \\	
	 \hline
  	Programa de Pós-Graduação em Matemática & Matemática &  & Doutor & DMAp\\
		Programa de Pós-Graduação em Matemática & Matemática &  & Mestre & MMAp\\
    \hline

\end{tabular}
\begin{flushleft}
		Fonte: Elaborado pelos autores.\
\end{flushleft}
\end{quadro}

% ----------------------------------------------------------
\chapter{Siglas dos Programas de Pós-Graduação do IFSC}
\index{quadros}O \autoref{quadro-ifsc} relaciona as siglas estabelecidas para os programas de pós-graduação do IFSC.
\begin{quadro}[htb]
\ABNTEXfontereduzida
\caption[Siglas dos Programas de Pós-Graduação do IFSC]{Siglas dos Programas de Pós-Graduação do IFSC}
\label{quadro-ifsc}
\begin{tabular}{|p{3.5cm}|p{3.5cm}|p{3.5cm}|p{1.5cm}|p{2.25cm}|}
  \hline
   \textbf{PROGRAMA} & \textbf{ÁREA DE CONCENTRAÇÃO} & \textbf{OPÇÃO} & \textbf{TÍTULO} & \textbf{SIGLA}  \\
    \hline
Graduate Program in Physics & Applied Physics & Biomolecular Physics & Doutor & DFAFBe\\
Programa de Pós-Graduação do Instituto de Física de São Carlos & Física Aplicada &  & Doutor & DFA\\
Programa de Pós-Graduação do Instituto de Física de São Carlos & Física Aplicada & Física Computacional & Doutor & DFAFC\\
Programa de Pós-Graduação do Instituto de Física de São Carlos & Física Aplicada & Física Biomolecular & Doutor & DFAFBp\\
Programa de Pós-Graduação do Instituto de Física de São Carlos & Física Aplicada &  & Mestre & MFA\\
Programa de Pós-Graduação do Instituto de Física de São Carlos & Física Aplicada & Física Computacional & Mestre & MFAFC\\
Programa de Pós-Graduação do Instituto de Física de São Carlos & Física Aplicada & Física Biomolecular & Mestre & MFAFB\\
Programa de Pós-Graduação do Instituto de Física de São Carlos & Física Básica &  & Doutor & DFB\\
Programa de Pós-Graduação do Instituto de Física de São Carlos & Física Básica &  & Mestre & MFB\\
		\hline

\end{tabular}
\begin{flushleft}
		Fonte: Elaborado pelos autores.\
\end{flushleft}
\end{quadro}

% ----------------------------------------------------------
\chapter{Siglas dos Programas de Pós-Graduação do IQSC}
\index{quadros}O \autoref{quadro-iqsc} relaciona as siglas estabelecidas para os programas de pós-graduação do IQSC.
\begin{quadro}[htb]
\ABNTEXfontereduzida
\caption[Siglas dos Programas de Pós-Graduação do IQSC]{Siglas dos Programas de Pós-Graduação do IQSC}
\label{quadro-iqsc}
\begin{tabular}{|p{3.5cm}|p{3.5cm}|p{3.5cm}|p{1.5cm}|p{2.25cm}|}
  \hline
   \textbf{PROGRAMA} & \textbf{ÁREA DE CONCENTRAÇÃO} & \textbf{OPÇÃO} & \textbf{TÍTULO} & \textbf{SIGLA}  \\
    \hline
Programa de Pós-Graduação do Instituto de Química de São Carlos & Físico-química &  & Doutor & DFQ\\
Programa de Pós-Graduação do Instituto de Química de São Carlos & Físico-química &  & Mestre & MFQ\\
Programa de Pós-Graduação do Instituto de Química de São Carlos & Química Analítica e Inirgânica &  & Doutor & DQAI\\
Programa de Pós-Graduação do Instituto de Química de São Carlos & Química Analítica e Inirgânica &  & Mestre & MQAI\\
Programa de Pós-Graduação do Instituto de Química de São Carlos & Química Orgânica e Biológica &  & Doutor & DQOB\\
Programa de Pós-Graduação do Instituto de Química de São Carlos & Química Orgânica e Biológica &  & Mestre & MQOB\\
\hline

\end{tabular}
\begin{flushleft}
		Fonte: Elaborado pelos autores.\
\end{flushleft}
\end{quadro}


% ----------------------------------------------------------
\chapter{Siglas dos Cursos de Graduação da EESC}
\index{quadros}O \autoref{quadro-geesc} relaciona as siglas estabelecidas para os cursos de graduação da EESC.
\begin{quadro}[htb]
	\ABNTEXfontereduzida
	\caption[Siglas dos Cursos de Graduação da EESC]{Siglas dos Cursos de Graduação da EESC}
	\label{quadro-geesc}
	\begin{tabular}{|p{6.5cm}|p{6.5cm}|p{1.75cm}|}
		\hline
		\textbf{CURSO} & \textbf{TÍTULO} &  \textbf{SIGLA}  \\
		\hline
		Engenharia Ambiental & Engenheiro Ambiental & EAMB\\
		Engenharia Aeronáutica & Engenheiro Aeronáutico & EAER\\
		Engenharia Civil & Engenheiro Civil & ECIV\\
		Engenharia de Computação & Engenheiro de Computação & ECOM\\
	    Engenharia Elétrica com Ênfase em Eletrônica & Engenheiro Eletricista & EELT\\
	    Engenharia Elétrica com Ênfase em Sistemas de Energia e Automação & Engenheiro Eletricista & EELS\\
		Engenharia de Materiais e Manufatura & Engenheiro de Materiais e de Manufatura & EMAT\\
		Engenharia Mecânica & Engenheiro Mecatrônico & EMET\\
		Engenharia de Produção & Engenheiro de Produção & EPRO\\
		\hline
		
	\end{tabular}
	\begin{flushleft}
		Fonte: Elaborado pelos autores.\
	\end{flushleft}
\end{quadro}

% ----------------------------------------------------------
\chapter{Exemplo de tabela centralizada verticalmente e horizontalmente}
\index{tabelas}A \autoref{tab-centralizada} exemplifica como proceder para obter uma tabela centralizada verticalmente e horizontalmente.
% utilize \usepackage{array} no PREAMBULO (ver em USPSC-modelo.tex) obter uma tabela centralizada verticalmente e horizontalmente
\begin{table}[htb]
\ABNTEXfontereduzida
\caption[Exemplo de tabela centralizada verticalmente e horizontalmente]{Exemplo de tabela centralizada verticalmente e horizontalmente}
\label{tab-centralizada}

\begin{tabular}{ >{\centering\arraybackslash}m{6cm}  >{\centering\arraybackslash}m{6cm} }
\hline
 \centering \textbf{Coluna A} & \textbf{Coluna B}\\
\hline
  Coluna A, Linha 1 & Este é um texto bem maior para exemplificar como é centralizado verticalmente e horizontalmente na tabela. Segundo parágrafo para verificar como fica na tabela\\
  Quando o texto da coluna A, linha 2 é bem maior do que o das demais colunas  & Coluna B, linha 2\\
\hline
\end{tabular}
\begin{flushleft}
		Fonte: Elaborada pelos autores.\
\end{flushleft}
\end{table}

% ----------------------------------------------------------
\chapter{Exemplo de tabela com grade}
\index{tabelas}A \autoref{tab-grade} exemplifica a inclusão de traços estruturadores de conteúdo para melhor compreensão do conteúdo da tabela, em conformidade com as normas de apresentação tabular do IBGE.
% utilize \usepackage{array} no PREAMBULO (ver em USPSC-modelo.tex) obter uma tabela centralizada verticalmente e horizontalmente
\begin{table}[htb]
\ABNTEXfontereduzida
\caption[Exemplo de tabelas com grade]{Exemplo de tabelas com grade}
\label{tab-grade}
\begin{tabular}{ >{\centering\arraybackslash}m{8cm} | >{\centering\arraybackslash}m{6cm} }
\hline
 \centering \textbf{Coluna A} & \textbf{Coluna B}\\
\hline
  A1 & B1\\
\hline
  A2 & B2\\
\hline
  A3 & B3\\
\hline
  A4 & B4\\
\hline
\end{tabular}
\begin{flushleft}
		Fonte: Elaborada pelos autores.\
\end{flushleft}
\end{table}


\end{apendicesenv}
% ---

% ----------------------------------------------------------
% Anexos
% ----------------------------------------------------------
\include{USPSC-Anexos}

%---------------------------------------------------------------------
% INDICE REMISSIVO
%--------------------------------------------------------------------
\phantompart
\printindex

%---------------------------------------------------------------------

\end{document}
